\documentclass[10pt,a4paper]{amsart}
\usepackage[margin=19mm]{geometry}
\usepackage{amsmath,amssymb,mathtools}
\usepackage[scr=rsfs,cal=euler]{mathalfa}
\usepackage{xfrac}
\usepackage[unicode,hidelinks,bookmarks=false]{hyperref}
\newcommand{\ie}{i.e.\ }
\newcommand{\cf}{cf.\ }
\newcommand{\resp}{respectively\ }
\newcommand{\Subsection}{\subsection}
\providecommand{\nobreakdash}{\nobreak\hbox{-}}
\setlength{\emergencystretch}{2em}
\multlinegap0pt
\usepackage{amssymb}
\newtheorem{theorem}{Theorem}
\def\CC{\mathbb{C}}
\def\ideal#1{\langle #1\rangle}
\title{An uncertainty principle for M\"obius inversion on posets}
\author{Marcel K. Goh}
\date{Source: arXiv:2302.02466v3 (CC BY 4.0)}
\begin{document}
\maketitle

\noindent\emph{Source and licence.} An uncertainty principle for M\"obius inversion on posets. Version arXiv:2302.02466v3 (CC BY 4.0), distributed by arXiv under the Creative Commons Attribution 4.0 International licence, http://creativecommons.org/licenses/by/4.0/, as recorded in the arXiv licence metadata for this exact version and shown on the abstract page https://arxiv.org/abs/2302.02466v3. Full licence notice: \texttt{licenses/arxiv-2302.02466-CC-BY-4.0.txt}.\par\smallskip

\noindent\textit{Editorial scope.} Counted unit: the article's theorem main (source0.tex lines 222--235). The article's complete original proof of it is delivered with it (source0.tex lines 237--255, one \texttt{proof} environment of 19 lines). No mathematics is written, repaired or re-derived: the statement and the proof are the article's own lines, in the article's own notation, and no retained line is substituted or re-flowed.  No further source-local context is needed: every cross-reference of every retained line resolves inside the two delivered spans.  Nothing was omitted from any retained span, so the package carries no omission record.  No retained line contains a citation, so the wrapper carries no bibliography and no bibliography entry.  No retained line contains a figure environment, an image-inclusion command or drawing code, so no image is delivered and the package declares no asset dependency.  Editorial formatting only, in addition: an AMS \texttt{amsart} preamble that restates the article's own declarations and macros used by the retained lines, namely the environments \texttt{theorem} on separate counters, together with the standard packages those lines need.  The retained environments print in this document as Theorem 1 in order of appearance, whereas the article prints the same items with the numbering of its own full document.  The complete native source of the article as distributed in the arXiv e-print for this exact version is bundled unchanged as \texttt{sources/137005/source0.tex}.
\begin{theorem}\label{main}
Let $P$ be a locally finite
join-semilattice with bottom element. Suppose that for every $y\in P$,
there exist infinitely
many $z\in P$ such that
\medskip
\item{i)} $z$ covers $y$; and
\smallskip
\item{ii)} for all $x\le y$, $\mu_P(x,z) = \mu_P(x,y)\mu_P(y,z) = -\mu_P(x,y)$.

\medskip
\noindent Then for any $f:P\to \CC$ that is not identically zero,
the support of $f$ and the support of its M\"obius transform $g$ cannot both be finite.
\end{theorem}

\begin{proof}
If $g$ has infinite support we are done, so suppose that $g$ has finite support
and choose $y$ with $f(y)\ne 0$. Let $S$ be the support of $g$ and let $T = S\cup\{y\}$. Let $m = \bigvee T$
be the join of these elements. Let $z$ be an element of $P$ that covers $y$ and note that if
$\ideal z\setminus \ideal y$ contains an element $x\in S$, then $z = x\vee y$, so $z\in \ideal m$.
We have shown that if $z\notin \ideal m$, then $S\cap \ideal z \setminus \ideal y = \emptyset$.

Since $P$ is locally finite with a bottom element, $\ideal m$ is finite; thus
by our hypothesis there exist infinitely many $z$ satisfying both conditions (i) and (ii)
that lie outside of $\ideal m$. Fix an arbitrary such $z$ and expand
\begin{align*}
f(z) &= \sum_{x \in \ideal z} \mu_P(x,z) g(x) \cr
&= \sum_{x\in \ideal y}\mu_P(x,z)g(x)
+ \sum_{x\in \ideal z \setminus \ideal y} \mu_P(x,z) g(x).\cr
\end{align*}
By the observation just proved, the second summation is zero, and by condition (ii),
$$f(z) = \sum_{x\in \ideal y} \mu_P(x,y)\mu_P(y,z) g(x) = \mu_P(y,z) f(y) = -f(y) \ne 0.$$
Hence we see that $f(z) \ne 0$ for infinitely many $z$.
\end{proof}

\end{document}
